\chapter*{前言：这本书想带你走的路}
\addcontentsline{toc}{chapter}{前言：这本书想带你走的路}

如果你手里有一台望远镜，它每晚记录下来的其实不是"一张星空照片"，而是一长串
带着时间、位置、颜色和偏振标签的"光子点击"。这本书想回答的问题只有一个：\emph{我
们怎样从这一串点击里，读出恒星的大小、黑洞边缘的样子、乃至暗物质留下的蛛丝马迹？}

这本书是为\textbf{学过大学普通物理，但没学过量子光学}的读者写的。你只需要
记得电磁波、一点量子力学、一点概率——剩下的相干函数、光子统计、量子估计、强度干
涉，我们都会从零开始搭。我的目标很具体：\textbf{每一个概念都讲到你脑子里能浮现出画
面，每一步公式都推到你不必自己补任何一步。}

为此，这本书特别下了三处功夫。第一，前面专门加了一整个"预备"部分，
把复振幅、傅里叶与带宽、泊松噪声、量子谐振子，一直到\textbf{光是怎样被量子化成光
子}的，都补齐；后面所有出现 $\hat a$、$|\alpha\rangle$、$g^{(2)}$ 的地方，你都能
追溯到它第一次被"造出来"的那一页。第二，每一条重要公式旁边都有一张"公式解读"小
卡片，逐个符号告诉你它是谁；正文再把它的单位、量纲、物理意义和数量级讲一遍。第三，
尽量在抽象公式后面代入真实的恒星和真实的望远镜，让你始终感到自己在和真实的天空
打交道。

你可以按顺序读。最前面的预备部分是全书的地基，看起来"太基础"，但正是它决定了后
面读起来是顺畅还是卡壳。如果你已经很熟悉某一块，跳过去也没关系，需要时再回来查。

写这本书，也是我自己学习量子天文学的方式。它还在生长，因为我也还在学。

\begin{flushright}
王瑜\\
2026 年 7 月
\end{flushright}
